\taub{} évalue des agents conversationnels à base de LLM. L'agent dispose d'outils et doit respecter une politique métier pour traiter les demandes d'un utilisateur, lui-même simulé par un LLM. Ce premier rapport décrit le benchmark (architecture de la simulation, domaines, tâches, outils, mécanisme de notation), puis analyse ses biais et ses limites. L'analyse s'appuie sur le code et les données du dépôt officiel, figés à un commit précis, et sur nos premiers essais.

Nous montrons que le score \texttt{airline} agrège deux compétences distinctes : 24 tâches sur 50 attendent que l'agent ne modifie rien (le plus souvent, refuser une demande contraire à la politique), et 26 attendent qu'il agisse correctement. Les résultats publiés ne distinguent pas ces deux familles. La notion de succès diffère aussi selon les domaines, contrairement à ce qu'indique la documentation. Les règles métier ne sont vérifiées que par leurs effets sur la base de données, et plus d'un tiers des tâches \texttt{airline} et \texttt{retail} ont dû être corrigées. Nos essais font apparaître un faux négatif (un agent pénalisé pour un champ facultatif) et un faux positif (une violation de la politique notée comme un succès). Nous menons enfin un premier essai sur l'effet d'un budget de raisonnement : \texttt{Qwen3-14B} en local, sur les 50 tâches \texttt{airline} et 3 essais par tâche. Le budget de 4096 tokens fait passer $\passk{1}$ de 30,7\,\% à 32,7\,\% et $\passk{2}$ de 18,0\,\% à 21,3\,\%. Ces écarts vont dans le sens attendu mais ne sont pas significatifs, et un score global ne dit pas sur quelle famille de tâches porte le gain. Nous en tirons des questions de recherche et un protocole pour la suite du projet.
